Definiamo un {\em monitor} come un oggetto in cui pu\`o {\em entrare}
un solo processo alla volta, e a cui \`e associata una {\em
variabile di condizione}. 
Il processo che si trova all'interno del monitor \`e detto {\em possessore}
del monitor.
Il possessore del monitor pu\`o eseguire le
operazioni di {\em attesa} e {\em notifica} sulla variabile di condizione
associata. L'operazione di attesa rilascia il monitor e sospende il processo fino a quando un
altro processo esegue l'operazione di notifica sulla stessa variabile di
condizione. L'operazione di notifica risveglia uno dei processi in attesa
sulla variabile di condizione, se ve ne sono, altrimenti non fa niente.
Se un processo ne risveglia un altro tramite una operazione di notifica, gli
cede anche il possesso del monitor, quindi si sospende su una coda di processi
``urgenti'' associata al monitor. I processi ``urgenti'' hanno precedenza nel
riottenere il possesso del monitor, non appena questo si libera.

Per realizzare i monitor definiamo la seguente struttura (file {\tt
sistema.cpp}):

\begin{verbatim}
    struct des_monitor {
        natl owner;
        des_proc *w_enter;
        des_proc *w_cond;
        des_proc *urgent;
    };
\end{verbatim}

Il campo {\tt owner} contiene l'id del processo che possiede il monitor (0 se
nessun processo possiede attualmente il monitor).  
Il campo \verb|w_cond| rappresenta la lista dei processi in attesa
sulla variabile di condizione associata al monitor.
Il campo \verb|w_enter| rappresenta la lista dei processi in attesa
di poter entrare nel monitor. Il campo {\tt urgent} rappresenta
la lista dei processi ``urgenti''.

Le seguenti primitive, accessibili dal livello utente, operano sui
monitor (nei casi di errore, abortiscono il processo chiamante):
\begin{itemize}
\item \verb|natl monitor_ini()| (gi\`a realizzata): inizializza un nuovo monitor, 
con tutti i campi a 0.  Restituisce l'identificatore del nuovo monitor. Se non
\`e possibile creare un nuovo monitor, restituisce {\tt 0xFFFFFFFF}.
\item \verb|void monitor_enter(natl mon)| (gi\`a realizzata): tenta di entrare nel monitor di identificatore
{\tt mon}.  Se il monitor appartiene gi\`a a qualche altro processo,
sospende il processo in attesa di potervi entrare.
\`E un errore tentare di entrare in un monitor che si possiede gi\`a. 
\item \verb|void monitor_leave(natl mon)|:
esce dal monitor di identificatore {\tt mon}.
\`E un errore tentare di uscire da un monitor che non si possiede.
La primitiva deve cedere il possesso del monitor ad uno degli eventuali
processi ``urgenti'' o in attesa di entrare nel monitor, se ve ne sono.
\item \verb|void monitor_wait(natl mon)|: si pone in
attesa sulla variabile di condizione associata al monitor di identificatore
{\tt mon}.  \`E un errore tentare di porsi in attesa su una variable di
condizione associata ad un monitor che non si possiede.
Prima di sospendersi, il processo deve cedere il possesso del monitor
ad uno dei processi ``urgenti'' o in attesa di entrare nel monitor, se ve ne sono.
\item \verb|void monitor_notify(natl mon)|(gi\`a realizzata):
esegue una notifica sulla
variabile di condizione associata al monitor di identificatore {\tt mon}.
\`E un errore tentare di eseguire una notifica su una variable di
condizione associata ad un monitor che non si possiede.
\end{itemize}

Modificare il file \verb|sistema.cpp| in modo da
realizzare le primitive mancanti.
